\begin{abstract}
Humans learn an unfamiliar physical world by acting in it, and then see new scenes in the light of what they learned. We ask whether multimodal agents can learn the physical world through interaction in the same way. We introduce a benchmark built around a single unit, the visual experiment: a rendered scene, an intervention, and a time interval, followed by the rendered outcome. During interaction, the agent chooses experiments in a world whose force law departs from canonical physics, either through a structural change to the law or through a hidden coefficient field, and sees each outcome as a rendered frame; checkpoint questions trace an online learning curve. After interaction is frozen, the roles reverse: the world shows new scenes, and the agent must understand each scene and predict the next frame of the world. Sibling worlds share the same scenes but follow different hidden laws, so a correct prediction requires knowing the particular world the agent explored. Evidence arms on identical questions separate what interaction contributes: no history, the agent's own history, a history from a sibling world, the same history shown to a passive observer, and the law stated in text. For the strongest system we evaluate, interaction raises accuracy from \V{GptSolNoneAcc}\% to \V{GptSolOwnAcc}\%, and a sibling world's history lowers it to \V{GptSolSibAcc}\%, often to the sibling's own answer, so the gain is specific to the world the agent explored. For both systems with complete runs, interaction recovers less than half of what knowing the law provides: told the law in text, they reach \V{OpusLawAcc}--\V{GptSolLawAcc}\%, so they can apply a law to new scenes but do not form it from their own experiments. A passive observer of the same experiments learns as much as the agent that ran them.
\end{abstract}
